%!TEX root = ../main.tex
\section{Conclusion}
\label{sec:conclusion}

Episode counts do not determine what an embodied model learns.
Learning depends on how raw records are made usable, converted into training units, expanded, and assigned to training stages.
This survey organizes those decisions by the operation performed, the signal that guides it, the evidence used to validate it, and the point in the lifecycle at which it changes.
The four axes connect source design and alignment to curation, generation, post-training, and repeated data updates without treating them as unrelated engineering steps.

The literature supports three conclusions that rule out common shortcuts.
Data utility is conditional on the learner, robot, task, stage, and surrounding dataset, so one universal quality score is inadequate.
The useful decision often lies below the episode because valid behavior, mistakes, retries, and recoveries require different labels and objectives.
Generated data must pass checks matched to its training role.
Visual plausibility alone does not establish valid action supervision.

Public industrial recipes reinforce these conclusions despite their different scales and disclosure levels.
They consistently separate source roles, align motion and control, validate data before action learning, and change sampling across pretraining, specialization, and post-training.
Feedback-driven systems extend this recipe with states visited by the current policy.
Those states can make failure correction data-efficient, but they also couple the next dataset to the policy and verifier's current errors.
Reference data, retained-capability tests, mixture controls, and joint data--model rollback are therefore part of the learning system rather than optional record keeping.

The next research step is not unrestricted automation.
It is a controlled data-preparation system whose decisions can be compared, attributed, and reversed.
Progress requires shared raw pools, fixed-budget learners, layered evidence, cost and risk accounting, and multi-round tests that expose error amplification.
An autonomous preparation agent becomes credible only when it can choose an informative experiment, respect authority limits, and decline to change the data when the available evidence is insufficient.
